\subsection{Complements on \(p\)-morphisms}\label{sga3:XXI:7.5}

Let \(f:\cR\to\cR'\) be a \(p\)-morphism
(cf.\ \ref{sga3:XXI:6.8}).  It is clear from the definitions that
the bijection \(u:R\xrightarrow{\sim}R'\) associated with \(f\)
makes the systems of simple roots, the systems of positive roots,
the irreducible components, and so forth of \(R\) and \(R'\)
correspond.  We may therefore suppose, for simplicity, that
\(R\) and \(R'\) are irreducible.

\begin{lemma}\label{sga3:XXI:7.5.1}
If \(R\) and \(R'\) are irreducible, there is a \(k\in\mathbb Q\)
such that, for every \(\alpha\in R\),
\[
 k\,\operatorname{long}(u(\alpha))
 =
 q(\alpha)^2\operatorname{long}(\alpha).
\]
\end{lemma}

\begin{proof}
One has
\[
 \operatorname{long}(\alpha)(\alpha^\ast,\beta)
 =
 \operatorname{long}(\beta)(\beta^\ast,\alpha)
\]
and, similarly,
\[
 \operatorname{long}(u(\alpha))
   (u(\alpha)^\ast,u(\beta))
 =
 \operatorname{long}(u(\beta))
   (u(\beta)^\ast,u(\alpha)).
\]
It follows immediately that, for nonorthogonal \(\alpha,\beta\),
\[
 \frac{q(\alpha)^2\operatorname{long}(\alpha)}
      {\operatorname{long}(u(\alpha))}
 =
 \frac{q(\beta)^2\operatorname{long}(\beta)}
      {\operatorname{long}(u(\beta))}.
\]
The conclusion then follows from \ref{sga3:XXI:7.1.3}\,\SGARefLink{sga3:XXI:7.1.3:item:iv}{(iv)}.
\end{proof}

\begin{remark}\label{sga3:XXI:7.5.2}
It follows from \ref{sga3:XXI:7.2.2} and
\ref{sga3:XXI:6.8.4} that \(q(\alpha)\) depends only on
\(\operatorname{long}(\alpha)\).  One then sees readily that, if
\(q(\alpha)\) is not constant, the product
\(q(\alpha)\operatorname{long}(\alpha)\) is constant; hence
\(p=2\) or \(3\).  Inspection of the preceding diagrams gives the
four possible cases (the same letter denotes a Dynkin diagram and
the corresponding reduced simply connected root datum):
\[
\begin{array}{llll}
p=2, & B_n\xrightarrow{\,f_1\,}C_n,
     & C_n\xrightarrow{\,f_2\,}B_n,
     & (C_2=B_2),\\
p=2, & F_4\xrightarrow{\,g\,}F_4,\\
p=3, & G_2\xrightarrow{\,h\,}G_2.
\end{array}
\]
The reader will observe that
\(f_1\circ f_2\), \(f_2\circ f_1\), \(g\circ g\), and
\(h\circ h\) are \(p\)-morphisms of the form described in
\ref{sga3:XXI:6.8.6}.
\end{remark}

\medskip
\noindent\textbf{7.5.3.}\label{sga3:XXI:7.5.3}
The preceding description shows immediately that, if \(\cR\) and
\(\cR'\) are two reduced root data of semisimple rank at most \(2\)
and there is a \(p\)-morphism from \(\cR'\) to \(\cR\), then
\(\cR\) and \(\cR'\) have the same type.  More precisely, one has
the following table.

\medskip
\noindent\emph{Notation.}
Let \(f:\cR'\to\cR\) be a \(p\)-morphism.  The symbols \(q\) and
\(q_1\) denote arbitrary positive powers of \(p\).  For rank-two
systems we use the notation of Section~\ref{sga3:XXI:4}, with
\(\alpha,\beta\) the simple roots and
\(\ell(\alpha)\leq\ell(\beta)\).

\begin{center}
\begingroup
\renewcommand{\arraystretch}{1.35}
\resizebox{\linewidth}{!}{%
\begin{tabular}{|c|c|c|c|}
\hline
\textbf{Type} & \(p\) & \textbf{values of \(f\)}
 & \textbf{values of \({}^tf\)}\\
\hline
Trivial & arbitrary & --- & ---\\
\hline
\(A_1\) & arbitrary
 & \(f(\alpha')=q\alpha\)
 & \({}^tf(\alpha^\ast)=q\alpha'^\ast\)\\
\hline
\(A_1\times A_1\) & arbitrary
 & \(\begin{array}{c}f(\alpha')=q\alpha\\
                      f(\beta')=q_1\beta\end{array}\)
 & \(\begin{array}{c}{}^tf(\alpha^\ast)=q\alpha'^\ast\\
                      {}^tf(\beta^\ast)=q_1\beta'^\ast\end{array}\)\\
\hline
\(A_2,B_2,G_2\) & arbitrary
 & \(\begin{array}{c}f(\alpha')=q\alpha\\
                      f(\beta')=q\beta\end{array}\)
 & \(\begin{array}{c}{}^tf(\alpha^\ast)=q\alpha'^\ast\\
                      {}^tf(\beta^\ast)=q\beta'^\ast\end{array}\)\\
\hline
\(B_2\) & \(p=2\)
 & \(\begin{array}{c}f(\alpha')=q\beta\\
                      f(\beta')=2q\alpha\end{array}\)
 & \(\begin{array}{c}{}^tf(\alpha^\ast)=2q\beta'^\ast\\
                      {}^tf(\beta^\ast)=q\alpha'^\ast\end{array}\)\\
\hline
\(G_2\) & \(p=3\)
 & \(\begin{array}{c}f(\alpha')=q\beta\\
                      f(\beta')=3q\alpha\end{array}\)
 & \(\begin{array}{c}{}^tf(\alpha^\ast)=3q\beta'^\ast\\
                      {}^tf(\beta^\ast)=q\alpha'^\ast\end{array}\)\\
\hline
\end{tabular}}
\endgroup
\end{center}

\section*{Bibliography}
\addcontentsline{toc}{section}{Bibliography}
\label{sga3:XXI:bibliography}

\begin{description}
\item[{[BLie]}] N.~Bourbaki,
  \emph{Groupes et algèbres de Lie}, Chapters IV--VI,
  Hermann, 1968.
\item[{[De80]}] M.~Demazure,
  ``A, B, C, D, E, F, etc.,'' pp.~221--227 in
  \emph{Séminaire sur les singularités des surfaces}
  (Palaiseau, 1976--1977), eds.\ M.~Demazure,
  H.~C.~Pinkham, and B.~Teissier,
  Lecture Notes in Mathematics \textbf{777},
  Springer-Verlag, 1980.
\end{description}
